\section{Discussion and open problems}\label{sec:discussion}

\subsection{Where the exponent comes from}\label{sec:disc-kappa}

The asymptotic saving $\kappa$ is governed by fundamentally different mechanisms depending on how record-address bits are grouped into chunks (Section~\ref{sec:main}). In~\cite{OAI26} and on the Bluestein comparison route (Appendix~\ref{app:bluestein}), the layout map $L$ is executed by width-$K$ chunk interchanges. Under that layout, the layer constraint $c<\beta(1-\tau)/\tau$ and the guard constraint $\eps<1/C_1$ force the chunk-exchange gap $\eps c(1-\tau)$ below $\beta(1-\tau)^2/(C_1\tau)$, capping $\kappa$ quadratically in the bit network's deficit $1-\tau$. Executing $L$ and $(LB)^{-1}$ via exact Walsh-layer positional permutations (Theorem~\ref{thm:positional-perm}) eliminates this external chunk-exchange cost altogether. The binding constraint becomes the simultaneous butterfly-layer gap $\eps(1-\lambda')$: since Proposition~\ref{prop:prop14-param} requires $\lambda'>\lambda>\sigma$ and $\eps<1/C_1$, every admissible choice satisfies
\[
  \kappa < \eps(1-\lambda') < \frac{1-\sigma}{C_1},
\]
which is \emph{linear} in the complex network's deficit $1-\sigma$.

This inequality bounds $\kappa$ in terms of the \emph{selected} rational exponent~$\sigma$; it is not yet the exact ceiling of the $h=100$ complex network itself. For our certified dyadic choice $1-\sigma=2^{-42}$ and guard exponent $C_1=14$, $(1-\sigma)/C_1\approx 2^{-45.8}$, so $\kappa=2^{-46}$ comes within a factor below~$2$ of the dyadic limit (and at $C_1=20$, $(1-\sigma)/C_1\approx 2^{-46.3}$, approached by $\kappa=2^{-47}$). The $h=100$ complex network admits any $\sigma$ with $s_c/W_c<m^\sigma$, that is, $1-\sigma<\gamma_c=-\ln(1-\eta_c)/\ln m\approx 2.654\cdot 10^{-13}$ (compared with $2^{-42}\approx 2.274\cdot 10^{-13}$). Already the rational test of Proposition~\ref{prop:deficits} at $L=14$ certifies $\sigma=1-\frac98\cdot 2^{-42}$, yielding non-dyadic exponents slightly above $2^{-42}/C_1$ up to the fixed-network ceiling $\gamma_c/C_1$. When the structural constants are instead fixed to the values $\beta=\frac12$, $1-\tau=1-\sigma=2^{-50}$, and $C_1=20$ printed in~\cite[Prop.~14]{OAI26}, the leaf condition $\lambda'>\sigma+\beta(1-\sigma)=1-2^{-51}$ imposes the tighter ceiling $\kappa<\eps(1-\beta)(1-\sigma)<2^{-51}/20<2^{-55}$, making $\kappa=2^{-56}$ the largest exponent of the form $2^{-j}$ admitted by our analysis with those constants. Finally, our guard exponent $C_1=14$ matches the threshold $11+3=14$ noted in~\cite[\S5.4]{OAI26} and proved in Proposition~\ref{prop:prop14-param}; we make no claim regarding smaller guard exponents.

\subsection{Activation sizes, scope, and practical limitations}\label{sec:disc-limitations}\label{sec:disc-perm}

Theorem~\ref{thm:main} is purely an asymptotic bit-complexity result and yields no practical speedup at any feasible input size. Even reaching dimension $d=\lfloor k^\eps\rfloor\ge 2$ at $\eps=1/15$ requires transform length $R=2^k$ with $k\ge 2^{15}$, while the chunk-width condition $K=\lfloor d^c\rfloor\ge 6$ at $\kappa=2^{-46}$ requires $\log_2 k\ge 3840\cdot 2^{42}\log_2 6$, before accounting for the guard constant $C_0\approx 10^{139}$ or the unspecified tape-stream thresholds of Facts~\ref{fact:oai-streams}--\ref{fact:oai-networks}. Moreover, even if all leading constants were~$1$, the asymptotic saving factor $(\log_2 n)^\kappa$ at $\kappa=2^{-46}$ would be smaller than $1+10^{-12}$ for all $n\le 10^{12}$. Accordingly, Algorithm~1 itself is executed only at test scale in exact arithmetic (Appendix~\ref{app:impl}).

By contrast, the C/GMP implementation (\texttt{algo=1}) of the two-dimensional~\cite{MS}-type $O(n\log n)$ multiplier uses a windowed sinc-Gaussian real kernel with additive precision schedule and Nyquist-column alternating-sum identity, and has modest enough constants to execute bit-exactly against GMP up to $10^9$ bits. That benchmarked program is not Algorithm~1 and not the negative-carrier or sub-$n\log n$ algorithm. On one Zen~3 core (one vCPU of a lightly loaded KVM guest; measurement details in~\cite{CGMP26}) at $10^8$ bits, it takes $92.5\,\mathrm{s}$ single-threaded vs.\ $0.606\,\mathrm{s}$ for GMP ($153\times$ excluding $5.5\,\mathrm{s}$ table setup, or $6.31\,\mathrm{s}$ and $10.4\times$ on $64$ threads). At $10^9$ bits, it takes $2613\,\mathrm{s}$ single-threaded vs.\ $8.10\,\mathrm{s}$ for GMP ($323\times$ excluding $26.0\,\mathrm{s}$ table setup), with $67.3\,\mathrm{GiB}$ (about $72\,\mathrm{GB}$) peak RAM~\cite{CGMP26} (at $10^9$ bits the natural threshold selects a second recursion level, $2^{27}\to 2^{14}\to 128$, whose parameters lie outside the regime covered by the written analysis of~\cite{MS}; the output is nevertheless bit-exact).

We also emphasize the scope of Theorem~\ref{thm:positional-perm}: it permutes the $m_0$ bit \emph{positions} of a complete binary record address carrying a long contiguous payload. It is not a sub-$n\log n$ routine for arbitrary permutations of $2^{m_0}$ records, nor does it replace the general field interchange of Lemma~\ref{lem:transfer12} used for the bridge's $O(d)$ axis exposures (where lines carry no long retained payload). Because the gadget forms complex Walsh--Hadamard linear combinations between exact permutation boundaries, it operates outside pure data-movement models~\cite[\S1.4]{HvdH25}.

\subsection{What became simpler, and what it cost}\label{sec:disc-simpler}

Relative to~\cite{HvdH21} and~\cite{OAI26}, the outer reduction discards every number-theoretic and polynomial-ring layer around the two $h=100$ finite networks (Table~\ref{tab:components}). Working on one power-of-two length $R=t_1\cdots t_d=2^k$ eliminates the Chinese remainder map onto distinct odd primes, trial-division prime search, and coprime-length Gaussian resampling with its Fourier-grid expansion. Modulating the sinc-Gaussian kernel by a negative half-bin phase eliminates the Nyquist special case of~\cite{MS} and gives one two-term cyclic solver with $\|B_a^{-1}\|_\infty\le\frac47 t_a$ uniformly across all cycle lengths. Direct axis exposures replace $O(d^2)$ adjacent swaps by $O(d)$ interchanges. Exact Walsh-layer positional permutations replace $O(k/K)$ chunk interchanges, Bluestein's chirp, last-axis twists, $\C[y]/(y^r+1)$ ring convolutions, $M$-fold error amplification, $\lg M$-bit output widening, and $\Theta(rP)$-bit packed products by a plain complex tensor FFT. That FFT delegates only $O(P^2)$-bit integer products under a single precision contract $P=O(p+\lg R)$. Correspondingly, the external dependency surface of the written proof shrinks to zero results of~\cite{HvdH21} and only the stream primitives and $h=100$ network certificates of~\cite[Lemmas~3--4, Props.~7--8]{OAI26} as black boxes (and to zero external results in the Lean~4 formalization).

These simplifications come with concrete costs of their own. The sub-$n\log n$ saving still originates entirely in the two $h=100$ finite networks of~\cite{OAI26}. Turning the simultaneous butterfly layer into an exact positional-permutation engine requires six completed layer calls per slice-batched $\CNOT$ stage (at most $36 N_{\mathrm{slice}}$ layer calls per positional map), together with a temporary precision lift of $8d$ fractional bits so intermediate dyadic denominators never truncate (Section~\ref{sec:walsh}). Reaching $\kappa=2^{-46}$ beyond the $\kappa=2^{-56}$ printed-constants bound also still requires the parameter-exposed re-derivation of the interchange, packed-addition, and butterfly-layer procedures in Appendix~\ref{app:parameter-transfer}. Once those network-backed tape primitives are in place, however, the entire reduction from integer multiplication to them consists of a fixed convolution kernel, a power-of-two FFT, and a two-term recurrence.

\subsection{Open problems}\label{sec:disc-open}

\begin{enumerate}[label=(\arabic*),leftmargin=*]
  \item \emph{Larger deficits and smaller guards.} Construct finite linear networks with relative rank deficit substantially larger than $\eta_c\approx 3.67\cdot 10^{-12}$, or establish butterfly-layer guard bounds with exponent $C_1<14$. By Section~\ref{sec:disc-kappa}, either advance translates linearly into a larger exponent~$\kappa$.
  \item \emph{Explicit numerical cutoff.} Extract explicit constants for all multi-stack or multitape stream primitives so as to give an explicit numerical upper bound on the threshold $n_*$ in Theorem~\ref{thm:main}.
  \item \emph{Other applications of Walsh-layer permutations.} Determine whether exact positional permutations via selective butterfly layers improve the exponents of other problems accelerated by the networks of~\cite{OAI26}, notably matrix transposition~\cite[\S11]{OAI26} and the exact discrete Fourier transform~\cite{OAI26b}.
  \item \emph{Downstream arithmetic primitives.} Decide whether the CRT-free fractional-shift bridge also simplifies downstream operations such as integer division and square roots~\cite[\S10]{OAI26}. Since these reduce to multiplication by Newton iteration (and $\mathsf M(n)/n$ is nondecreasing for our bound), they inherit the $O\bigl(n(\log n)^{1-2^{-46}}\bigr)$ complexity in the standard way, though we have not checked the details (nor those for base conversion).
\end{enumerate}
